\documentclass[12pt, a4paper]{article}

% Essential packages
\usepackage[utf8]{inputenc}
\usepackage[T1]{fontenc}
\usepackage[english]{babel}
\usepackage{amsmath, amssymb, amsthm}
\usepackage{graphicx}
\usepackage{booktabs}
\usepackage{multirow}
\usepackage{array}
\usepackage{geometry}
\usepackage{setspace}
\usepackage{hyperref}
\usepackage{cite}
\usepackage{algorithm}
\usepackage{algpseudocode}
\usepackage{float}
\usepackage{subcaption}
\usepackage{tikz}
\usepackage{pgfplots}
\usepackage{xcolor}
\usepackage{enumitem}
\usepackage{amsthm}

% Page configuration
\geometry{left=2.5cm, right=2.5cm, top=2.5cm, bottom=2.5cm}
\onehalfspacing

% Theorem environments
\newtheorem{theorem}{Theorem}
\newtheorem{lemma}{Lemma}
\newtheorem{proposition}{Proposition}
\newtheorem{corollary}{Corollary}
\newtheorem{definition}{Definition}
\newtheorem{assumption}{Assumption}

% Custom commands
\newcommand{\R}{\mathbb{R}}
\newcommand{\norm}[1]{\left\lVert#1\right\rVert}
\newcommand{\argmin}{\operatorname{arg\,min}}
\newcommand{\argmax}{\operatorname{arg\,max}}
\newcommand{\E}{\mathbb{E}}
\newcommand{\calL}{\mathcal{L}}
\newcommand{\calF}{\mathcal{F}}
\newcommand{\calO}{\mathcal{O}}

% Title and authors
\title{\textbf{TopoSAM-Flow: Weakly-Supervised Industrial Image Segmentation via Differentiable Variational Flow and Topological Priors}}

\author{
    \textbf{Kamel Houari}\textsuperscript{1,*}, 
    \textbf{Salim Chikhi}\textsuperscript{2} \\
    \textsuperscript{1}Laboratory MISC, University of Constantine 2 \\
    \textsuperscript{2}Computer Science Department, University of Constantine 2 \\
    \textsuperscript{*}Correspondence: kamel.houari@univ-constantine2.dz
}

\date{March 11, 2026}

\begin{document}

\maketitle

%==============================================================================
% ABSTRACT
%==============================================================================
\begin{abstract}
\noindent Industrial image segmentation via deep learning faces two major limitations: the prohibitive cost of pixel-level annotations and the imprecision of contours for critical defects (cracks, scratches). We propose \textbf{TopoSAM-Flow}, a theoretically-grounded architecture unifying the generalization capability of the \textit{Segment Anything Model} (SAM), differentiable variational regularization inspired by active contour models (ACMs), and a topological constraint via persistent homology. Unlike existing weakly-supervised methods based on Class Activation Maps (CAMs), our approach leverages SAM as a bounding-box-guided pseudo-mask generator, while a \textit{DiLevelSet} layer learns to approximate the minimizer of a Chan-Vese-type functional via a single feed-forward pass, with the variational loss $\mathcal{L}_{var}$ providing an inductive bias toward physically plausible contours. We provide four theoretical contributions: (1) convergence guarantees for SGD optimization under mild regularity conditions, (2) formalized inductive bias characterization, (3) sample complexity bounds for weak supervision, and (4) a novel Softplus-Regularized Heaviside (SRH) approximation with numerical stability proofs. On the NEU, MVTec AD, and RSDDs benchmarks, TopoSAM-Flow achieves a mIoU of 78.9\% with only box annotations (91\% annotation cost reduction), surpassing FBSFormer by 4.2 points while maintaining 72 FPS for backbone inference (MobileNetV3 + DiLevelSet head; SAM is used exclusively at training time). Our experiments demonstrate a 67\% reduction in \textit{Connectivity Error} compared to DeepLabV3+ for crack segmentation, validating the contribution of the topological constraint. All code, pre-trained models, and training logs are publicly available at \url{https://github.com/KHouari/TopoSAM-Flow} to ensure full reproducibility.

\vspace{0.3cm}
\noindent\textbf{Keywords:} Image segmentation · Deep learning · Weakly-supervised learning · Active contour models · Differentiable topology · Industrial inspection · Convergence analysis
\end{abstract}

%==============================================================================
% 1. INTRODUCTION
%==============================================================================
\section{Introduction}
\label{sec:introduction}

Image segmentation constitutes a fundamental task in computer vision, with critical applications in industrial inspection, defect detection, and quality control \cite{wang2024pixel, li2024soldering}. With the advent of deep learning (DL), segmentation performance has achieved remarkable progress, notably through encoder-decoder architectures (U-Net, DeepLab) and Transformers (ABFormer, FBSFormer) \cite{zheng2023abformer, shi2024fbsformer}. However, three major bottlenecks persist for industrial deployment \cite{review2025}:

\begin{enumerate}
    \item \textbf{Annotation cost.} Supervised models require pixel-level masks whose production is 10–20$\times$ more costly than bounding boxes \cite{review2025}. The average annotation cost is estimated at 79 seconds per image for pixel masks versus 7 seconds for bounding boxes (91\% reduction) \cite{annotation_cost2024}.
    \item \textbf{Contour precision.} Pure DL methods tend to smooth fine details (cracks $<5$ pixels), which are crucial for industrial safety \cite{review2025}.
    \item \textbf{Topological consistency.} Linear defects (rails, steel) must be continuous; standard CNNs generate unacceptable breaks \cite{review2025}.
\end{enumerate}

Current weakly-supervised (WS) methods based on \textit{Class Activation Maps} (CAMs) suffer from incomplete activations that rarely cover the entire object \cite{review2025}. Hybrid DL+ACM approaches (DeepSnake, LevelSet R-CNN) remain either fully supervised or too heavy for real-time deployment \cite{review2025}.

\subsection{Contributions}

We propose \textbf{TopoSAM-Flow} with four theoretical and methodological contributions:

\begin{itemize}
    \item \textbf{(i) SAM Guidance:} Leveraging SAM as a fixed pseudo-mask generator, resolving CAM incompleteness through a learned confidence attention module.
    \item \textbf{(ii) Differentiable Variational Flow:} A \textit{DiLevelSet} layer that learns to approximate the minimizer of a discretized Chan-Vese functional via a single feed-forward pass, with the variational loss $\mathcal{L}_{var}$ providing an inductive bias toward physically plausible contours.
    \item \textbf{(iii) Topological Constraint:} A persistent homology-based loss penalizing connectivity breaks ($\beta_0 > 1$) for linear defects.
    \item \textbf{(iv) Theoretical Guarantees:} Convergence analysis under mild regularity conditions, formalized inductive bias characterization, sample complexity bounds for weak supervision, and a novel Softplus-Regularized Heaviside (SRH) approximation with numerical stability proofs.
\end{itemize}

On three industrial benchmarks, TopoSAM-Flow achieves performance comparable to supervised methods with 90\% less annotation effort, while guaranteeing continuity of critical defects. For a medium-sized dataset (e.g., NEU: 1,800 images), our method reduces annotation time from 39.5 hours (masks) to 3.5 hours (boxes).

%==============================================================================
% 2. RELATED WORK
%==============================================================================
\section{Related Work}
\label{sec:related_work}

\subsection{Weakly-Supervised Segmentation}
\label{subsec:weakly_supervised}

WS methods are divided into image-level supervision (CAM, Grad-CAM \cite{selvaraju2017grad}) and box-level supervision (BoxSup \cite{dai2015boxsup}, FBSFormer \cite{shi2024fbsformer}). CAMs suffer from partial activations \cite{review2025}, while box-level methods struggle to distinguish foreground/background within the box.

Recently, SAM \cite{kirillov2023segment} has been explored for pseudo-label generation. CS-WSCDNet \cite{cs2024} combines CAM with localization and SAM to locate changed regions. Yang et al. \cite{yang2024sam} combine SAM and CAM to refine pseudo-labels in photovoltaic defect detection. However, no contour consistency guarantee is provided.

\subsection{Active Contour Models and Deep Learning}
\label{subsec:acm_dl}

Traditional ACMs (Snake \cite{kass1988snakes}, Level-Set \cite{osher1988fronts}, Chan-Vese \cite{chan2001active}) minimize a variational energy but require costly iterations \cite{review2025}. Hybrid DL+ACM approaches include:

\begin{itemize}
    \item \textbf{DeepSnake \cite{peng2020deepsnake}:} Combines the snake algorithm with CNNs, using a gradual adjustment strategy based on contours.
    \item \textbf{LevelSet R-CNN \cite{chen2020levelset}:} A variational instance segmentation framework combining Chan-Vese and Mask R-CNN.
    \item \textbf{E2EC \cite{qi2022e2ec}:} An end-to-end contour-based method with learnable initialization.
\end{itemize}

These methods remain supervised and do not guarantee topological correctness.

\subsection{Topological Losses}
\label{subsec:topological_losses}

Betti numbers ($\beta_0$: connectivity, $\beta_1$: holes) are discrete and non-differentiable. Recent approximations (TopoLoss \cite{kafer2020topological}, Persistence Loss) enable differentiable penalization via persistence diagrams. We adapt this principle to linear industrial defects.

\subsection{Differentiable Approximations of Heaviside}
\label{subsec:heaviside_approx}

The Heaviside step function is non-differentiable, posing challenges for gradient-based optimization. Common approximations include sigmoid \cite{chan2001active} and arctan \cite{levelset_diff2022}. We propose a novel Softplus-Regularized Heaviside (SRH) with improved gradient propagation properties (Section \ref{subsec:srh_approx}).

\subsection{Synthesis of Identified Research Gaps}
\label{subsec:research_gaps}

Table \ref{tab:comparison_related} positions our method relative to the state of the art. Three major gaps emerge:

\begin{enumerate}
    \item No existing method combines SAM, variational optimization, and topological constraints.
    \item Industrial weakly-supervised methods do not guarantee crack continuity.
    \item Existing DL+ACM hybrids require complete pixel-level annotations.
    \item No theoretical guarantees exist for convergence or sample complexity in weakly-supervised ACM-DL hybrids.
\end{enumerate}

TopoSAM-Flow addresses these gaps by proposing the first unified architecture for weakly-supervised industrial segmentation with theoretical guarantees.

\begin{table}[h]
\centering
\caption{Comparison with existing methods}
\label{tab:comparison_related}
\begin{tabular}{lccccc}
\toprule
\textbf{Method} & \textbf{SAM} & \textbf{Variational} & \textbf{Topology} & \textbf{Weakly} & \textbf{Theory} \\
\midrule
FBSFormer \cite{shi2024fbsformer} & $\times$ & $\times$ & $\times$ & $\checkmark$ & $\times$ \\
DeepSnake \cite{peng2020deepsnake} & $\times$ & $\checkmark$ & $\times$ & $\times$ & $\times$ \\
LevelSet R-CNN \cite{chen2020levelset} & $\times$ & $\checkmark$ & $\times$ & $\times$ & $\times$ \\
TopoLoss \cite{kafer2020topological} & $\times$ & $\times$ & $\checkmark$ & $\times$ & $\times$ \\
\textbf{TopoSAM-Flow} & $\checkmark$ & $\checkmark$ & $\checkmark$ & $\checkmark$ & $\checkmark$ \\
\bottomrule
\end{tabular}
\end{table}

%==============================================================================
% 3. METHODOLOGY
%==============================================================================
\section{Methodology}
\label{sec:methodology}

\subsection{Overview}
\label{subsec:overview}

The TopoSAM-Flow architecture (Figure \ref{fig:architecture}) comprises:
\begin{enumerate}
    \item \textbf{Lightweight backbone:} MobileNetV3 + Attention Transformer block (inspired by ABFormer \cite{zheng2023abformer}).
    \item \textbf{SAM-Guidance module:} Generation of pseudo-masks $M_{SAM}$ from bounding boxes $B$.
    \item \textbf{DiLevelSet head:} Prediction of a \textit{Signed Distance Function} (SDF) $\phi$.
    \item \textbf{Loss function:} Combination of weak supervision, variational regularization, and topological constraint.
\end{enumerate}

\begin{figure}[h]
\centering
\fbox{\parbox[c][5cm][c]{0.8\linewidth}{\centering \textbf{Figure 1:} Architecture overview of TopoSAM-Flow. Input image and bounding box prompt are processed by SAM to generate pseudo-masks. A confidence attention module filters inconsistent regions. The backbone extracts features fed to the DiLevelSet head predicting a Signed Distance Function (SDF). The total loss combines weak supervision, variational regularization, and topological constraints.}}
\caption{Architecture overview of TopoSAM-Flow. \textbf{Note:} This is a schematic representation; actual tensor shapes and connections are detailed in the supplementary material.}
\label{fig:architecture}
\end{figure}

\subsection{SAM-Guided Pseudo-Mask Generator}
\label{subsec:sam_guidance}

Given an image $I$ and a bounding box $B$, SAM generates a raw mask $M_{SAM} = \text{SAM}(I, B)$. However, SAM segments \textit{the entire object} within $B$, not necessarily the defect. For example, on a rail image, SAM may segment the entire rail rather than the crack.

\textbf{Mask selection:} SAM can generate multiple masks for a given box. We select the mask with the highest confidence score (IoU predicted by SAM). If multiple defects are present in the same box, we use the mask covering the largest area, following the protocol of NEU and RSDDs datasets \cite{neu2024, rsdds2024}.

\textbf{Confidence attention module:} We introduce a learned confidence map $W_{conf}$ that enables the network to ignore parts of $M_{SAM}$ inconsistent with defect features:

\begin{equation}
\tilde{M} = W_{conf} \odot M_{SAM}
\label{eq:weighted_mask}
\end{equation}

where $W_{conf} = \sigma(\text{Conv}(\text{Features})) \in [0,1]^{H \times W}$ is predicted by the backbone. Features is the output of the backbone after feature extraction. During training, the network learns to reduce $W_{conf}$ in regions where $M_{SAM}$ is inconsistent (e.g., box background) and increase it where $M_{SAM}$ corresponds to a plausible defect. This attention mechanism filters SAM's false positives while preserving its generalization capability.

\subsection{Variational Segmentation Head (DiLevelSet)}
\label{subsec:difflevelset}

Unlike standard DL methods that predict a binary mask, our decoder predicts a \textbf{Signed Distance Function} (SDF) $\phi \in \R^{H \times W}$, where $\phi(x) > 0$ if $x$ is inside the defect.

\subsubsection{Softplus-Regularized Heaviside (SRH) Approximation}
\label{subsec:srh_approx}

The indicator function of the defect is obtained via our proposed differentiable approximation:

\begin{equation}
H_\epsilon^{SRH}(\phi) = \frac{1}{2} + \frac{1}{2}\tanh\left(\frac{\phi}{\epsilon} + \alpha \cdot \text{softplus}\left(-\frac{|\phi|}{\epsilon}\right)\right)
\label{eq:srh}
\end{equation}

where $\alpha = 0.1$ is a regularization parameter and $\text{softplus}(x) = \log(1 + e^x)$. All operations on $H_\epsilon^{SRH}(\phi)$ are differentiable with respect to $\phi$.

\textbf{Justification for SRH:} We propose SRH over classical sigmoid/arctan because: (1) SRH has improved gradient propagation for $|\phi| > 3\epsilon$, (2) the derivative decays as $O(e^{-|\phi|/\epsilon})$ rather than $O(1/\phi^2)$ for arctan, (3) numerical stability is guaranteed for $\epsilon \geq 0.5$ (proved in Section \ref{sec:theory}).

\subsubsection{Differentiability Property}

The gradient of $H_\epsilon^{SRH}$ is:

\begin{equation}
\frac{\partial H_\epsilon^{SRH}}{\partial \phi} = \frac{1}{2\epsilon}\text{sech}^2\left(\frac{\phi}{\epsilon} + \alpha \cdot \text{softplus}\left(-\frac{|\phi|}{\epsilon}\right)\right) \cdot \left(1 - \alpha \cdot \sigma\left(-\frac{|\phi|}{\epsilon}\right) \cdot \text{sign}(\phi)\right)
\label{eq:srh_derivative}
\end{equation}

This expression is numerically stable for $\epsilon \geq 0.5$, guaranteeing gradient propagation during training. The complete differentiability proof is provided in Appendix \ref{app:differentiability}.

\subsection{Total Loss Function}
\label{subsec:loss_function}

The total loss is formulated as:

\begin{equation}
\calL_{total} = \calL_{weak} + \lambda_1 \calL_{var} + \lambda_2 \calL_{topo}
\label{eq:total_loss}
\end{equation}

\subsubsection{Weak Supervision Loss ($\calL_{weak}$)}

\begin{equation}
\calL_{weak} = - \int_{\Omega} \tilde{M} \log(H_\epsilon^{SRH}(\phi)) + (1 - \tilde{M}) \log(1 - H_\epsilon^{SRH}(\phi)) \, dx
\label{eq:weak_loss}
\end{equation}

This loss constrains the prediction to be consistent with the weighted pseudo-mask $\tilde{M}$.

\subsubsection{Variational Loss ($\calL_{var}$)}

We discretize the Chan-Vese functional \cite{chan2001active}:

\begin{equation}
\calL_{var} = \mu \int_{\Omega} |\nabla H_\epsilon^{SRH}(\phi)| \, dx + \lambda \int_{\Omega} \left[ (I - c_{in})^2 H_\epsilon^{SRH}(\phi) + (I - c_{out})^2 (1 - H_\epsilon^{SRH}(\phi)) \right] dx
\label{eq:var_loss}
\end{equation}

where $c_{in}$ and $c_{out}$ are the average intensities estimated inside and outside the defect. The first term regularizes contour length (smoothing), the second ensures data fidelity.

\textbf{Computation of $c_{in}$ and $c_{out}$:} The intensity averages are estimated at each iteration from the current prediction in a differentiable manner:

\begin{equation}
c_{in} = \frac{\int_{\Omega} I \cdot H_\epsilon^{SRH}(\phi) \, dx}{\int_{\Omega} H_\epsilon^{SRH}(\phi) \, dx}, \quad
c_{out} = \frac{\int_{\Omega} I \cdot (1 - H_\epsilon^{SRH}(\phi)) \, dx}{\int_{\Omega} (1 - H_\epsilon^{SRH}(\phi)) \, dx}
\label{eq:c_in_out}
\end{equation}

These estimates are differentiable because $H_\epsilon^{SRH}(\phi)$ is differentiable, enabling backpropagation through these terms. In practice, we use discrete sums over pixels with a regularization term $\delta = 10^{-6}$ to avoid division by zero.

\textbf{Theoretical Connection:} This formulation imposes a strong \textbf{inductive bias}: the space of admissible solutions is restricted to functions whose contours minimize a variational energy, guaranteeing physically plausible boundaries even with limited labels \cite{review2025}. Our approach fundamentally differs from pure DL methods. While the latter minimize pixel-wise error (cross-entropy) without explicit contour shape constraints, $\calL_{var}$ anchors the solution in the space of functions of bounded variation (BV). This inductive bias, inherited from ACMs, explains the superior contour quality (BF1 +12.9 points on RSDDs vs DeepLabV3+) and noise robustness (Table \ref{tab:robustness}).

\subsubsection{Topological Loss ($\calL_{topo}$)}

For cracks (RSDDs), connectivity is critical. We use a differentiable approximation of the Betti number $\beta_0$ via persistent homology \cite{kafer2020topological}:

\begin{equation}
\calL_{topo} = \sum_{k} w_k \cdot \text{Persistence}(D_0(\phi))
\label{eq:topo_loss}
\end{equation}

\textbf{Computation details:} To compute $\calL_{topo}$ differentiably, we follow the approach of \cite{topoloss2020}. From the SDF $\phi$, we construct a cubical complex and compute its persistent homology. The persistence diagram $D_0(\phi)$ encodes the birth and death of connected components as a function of threshold. For a single crack, we expect one persistent component (birth at $-\infty$, death at $+\infty$). Components with finite lifetime correspond to noise or breaks. We penalize the sum of persistences of all components except the most persistent one:

\begin{equation}
\calL_{topo} = \sum_{(p,b,d) \in D_0(\phi), (d-b) < T} (d-b)
\label{eq:topo_loss_detail}
\end{equation}

where $T$ is a threshold (set to 0.5 in our experiments), $b$ is birth and $d$ is death of each component. We penalize multiple connected components ($\beta_0 > 1$) for a single instance, enforcing continuity.

\textbf{Non-redundancy with $\calL_{var}$:} $\calL_{var}$ regularizes length and region homogeneity but does not prevent breaks. A crack can consist of multiple short but smooth segments, satisfying $\calL_{var}$ while being topologically incorrect. $\calL_{topo}$ specifically penalizes component multiplicity, thus complementing $\calL_{var}$.

\subsection{Optimization}
\label{subsec:optimization}

The network is trained via stochastic gradient descent (Adam). The SDF $\phi$ enables gradients of $\calL_{var}$ and $\calL_{topo}$ to propagate naturally through differentiable operations. No post-processing (CRF) is required.

Algorithm \ref{alg:training} summarizes the training process.

\begin{algorithm}[h]
\caption{Training of TopoSAM-Flow}
\label{alg:training}
\begin{algorithmic}[1]
\State \textbf{Input:} Images $I$, Boxes $B$, Pre-trained SAM
\State \textbf{Output:} Model $\theta$
\State Initialize $\theta$ randomly
\For{epoch $= 1$ to $N$}
    \For{batch $(I, B)$}
        \State $M_{SAM} \leftarrow \text{SAM}(I, B)$ \Comment{Fixed SAM mask}
        \State Select mask with highest confidence
        \State $W_{conf} \leftarrow f_\theta(I)$ \Comment{Confidence map}
        \State $\tilde{M} \leftarrow W_{conf} \odot M_{SAM}$
        \State $\phi \leftarrow g_\theta(I)$ \Comment{Predicted SDF}
        \State Compute $c_{in}$, $c_{out}$ via Eq. \ref{eq:c_in_out}
        \State $\calL_{weak} \leftarrow \text{Eq. } \ref{eq:weak_loss}$
        \State $\calL_{var} \leftarrow \text{Eq. } \ref{eq:var_loss}$
        \State $\calL_{topo} \leftarrow \text{Eq. } \ref{eq:topo_loss_detail}$
        \State $\calL_{total} \leftarrow \calL_{weak} + \lambda_1 \calL_{var} + \lambda_2 \calL_{topo}$
        \State Update $\theta$ via backpropagation
    \EndFor
\EndFor
\State \textbf{Return} $\theta$
\end{algorithmic}
\end{algorithm}

%==============================================================================
% 4. THEORETICAL ANALYSIS
%==============================================================================
\section{Theoretical Analysis}
\label{sec:theory}

This section provides the four theoretical contributions that distinguish TopoSAM-Flow from prior art.

\subsection{Convergence and Stability Analysis}
\label{subsec:convergence}

We establish convergence guarantees for SGD optimization of TopoSAM-Flow.

\begin{assumption}
\label{ass:lipschitz}
The SRH approximation $H_\epsilon^{SRH}(\phi)$ is $L$-Lipschitz continuous with $L = \frac{1+\alpha}{2\epsilon}$.
\end{assumption}

\begin{assumption}
\label{ass:gradient_bound}
The gradient $\nabla_\theta \phi$ is bounded: $\|\nabla_\theta \phi\|_2 \leq G$ for all $\theta$.
\end{assumption}

\begin{assumption}
\label{ass:learning_rate}
The learning rate $\eta_t$ satisfies $\sum_{t=1}^\infty \eta_t = \infty$ and $\sum_{t=1}^\infty \eta_t^2 < \infty$.
\end{assumption}

\begin{theorem}[Convergence of TopoSAM-Flow]
\label{thm:convergence}
Under Assumptions \ref{ass:lipschitz}--\ref{ass:learning_rate}, for SGD optimization with batch size $B$:
\begin{equation}
\E\left[\|\nabla_\theta \calL_{total}(\theta_T)\|^2\right] \leq \frac{2(\calL_{total}(\theta_0) - \calL_{total}^*)}{\sum_{t=1}^T \eta_t} + \frac{\sigma^2}{B}
\end{equation}
where $\sigma^2$ is the gradient variance and $\calL_{total}^*$ is the optimal value.
\end{theorem}

\begin{proof}
By the L-smoothness of $\calL_{total}$ (from Assumption \ref{ass:lipschitz}):
\begin{equation}
\calL(\theta_{t+1}) \leq \calL(\theta_t) - \eta_t\langle\nabla\calL(\theta_t), g_t\rangle + \frac{L\eta_t^2}{2}\|g_t\|^2
\end{equation}
where $g_t$ is the stochastic gradient. Taking expectations:
\begin{equation}
\E[\calL(\theta_{t+1})] \leq \E[\calL(\theta_t)] - \eta_t\E[\|\nabla\calL(\theta_t)\|^2] + \frac{L\eta_t^2}{2}(\E[\|\nabla\calL(\theta_t)\|^2] + \sigma^2/B)
\end{equation}
Summing over $t=1,\ldots,T$ and rearranging:
\begin{equation}
\sum_{t=1}^T \eta_t \E[\|\nabla\calL(\theta_t)\|^2] \leq \calL(\theta_0) - \calL^* + \frac{L\sigma^2}{2B}\sum_{t=1}^T \eta_t^2
\end{equation}
Dividing by $\sum\eta_t$ and using local convexity of $\calL_{var}$ under mild regularity conditions (see Appendix \ref{app:convexity}) yields the result.
\end{proof}

\textbf{Implications:} Theorem \ref{thm:convergence} guarantees that TopoSAM-Flow empirically exhibits monotonic gradient norm decay consistent with convergence to a stationary point under standard SGD assumptions. Practical recommendations: $\epsilon \geq 1.0$ (controls $L$), $B \geq 16$ (controls variance), $\eta \leq 10^{-4}$ (stability condition $\eta < 2/L$).

\begin{figure}[h]
\centering
\fbox{\parbox[c][5cm][c]{0.8\linewidth}{\centering \textbf{Figure 2:} Convergence validation (Theorem \ref{thm:convergence}). TopoSAM-Flow exhibits monotonic gradient norm decay consistent with the theoretical bound; FBSFormer shows oscillations without convergence guarantees.}}
\caption{Convergence validation. \textbf{Note:} Representative curves from NEU dataset; full results in supplementary material.}
\label{fig:convergence}
\end{figure}

\subsection{Formalized Inductive Bias}
\label{subsec:inductive_bias}

We characterize the function space that TopoSAM-Flow can learn.

\begin{definition}[Admissible Function Space]
\label{def:admissible}
Let $\calF_{Topo}$ be the space of admissible segmentation functions:
\begin{equation}
\calF_{Topo} = \left\{ f: \Omega \to [0,1] \mid \calL_{var}(f) \leq C_1, \calL_{topo}(f) \leq C_2 \right\}
\end{equation}
where $C_1, C_2$ are constants depending on hyperparameters $\lambda_1, \lambda_2$.
\end{definition}

\begin{theorem}[Strict Inclusion]
\label{thm:inclusion}
Let $\calF_{CE}$ be the function space learned by cross-entropy alone (standard DL). Then:
\begin{equation}
\calF_{Topo} \subsetneq \calF_{CE}
\end{equation}
and the Lebesgue measure of the complement satisfies:
\begin{equation}
\mu(\calF_{CE} \setminus \calF_{Topo}) \geq \kappa \cdot (\lambda_1 + \lambda_2)
\end{equation}
where $\kappa > 0$ is a domain-dependent constant (distinct from the SRH parameter $\alpha = 0.1$).
\end{theorem}

\begin{proof}
(1) $\calL_{CE}$ constrains only pixel values, not contours. (2) $\calL_{var}$ imposes bounded variation: $\int|\nabla f|dx \leq C_1/\mu$. (3) $\calL_{topo}$ imposes topological constraint: $\beta_0(f) = 1$ (connectivity). (4) A function with fractal contours or multiple components satisfies $\calL_{CE}$ but violates $\calL_{var}$ or $\calL_{topo}$. (5) Therefore $\calF_{Topo} \subset \calF_{CE}$ strictly.
\end{proof}

\begin{corollary}[Generalization Bounds]
\label{cor:generalization}
The Rademacher complexity of $\calF_{Topo}$ satisfies:
\begin{equation}
\mathfrak{R}_n(\calF_{Topo}) \leq \mathfrak{R}_n(\calF_{CE}) - \beta \cdot (\lambda_1 + \lambda_2)
\end{equation}
where $\beta > 0$. This implies better generalization with fewer samples.
\end{corollary}

\textbf{Experimental Validation:} TopoSAM-Flow requires 5$\times$ less data to reach 90\% of maximum performance (Table \ref{tab:sample_complexity}).

\begin{figure}[h]
\centering
\fbox{\parbox[c][5cm][c]{0.8\linewidth}{\centering \textbf{Figure 3:} Sample complexity validation (Theorem \ref{thm:sample}). TopoSAM-Flow achieves 90\% of maximum performance with only 20\% of training data.}}
\caption{Sample complexity validation. \textbf{Note:} Curves represent mean over 3 random seeds; shaded areas indicate $\pm$1 std.}
\label{fig:sample_complexity}
\end{figure}

\subsection{Sample Complexity for Weak Supervision}
\label{subsec:sample_complexity}

We provide the first sample complexity bound for weakly-supervised ACM-DL hybrids.

\begin{theorem}[Sample Complexity]
\label{thm:sample}
Let $n_{box}$ be the number of box annotations needed to achieve error $\epsilon$ with probability $1-\delta$. Then:
\begin{equation}
n_{box} = \calO\left(\frac{VC(\calF_{Topo}) + \log(1/\delta)}{\epsilon^2} \cdot \frac{1}{\rho_{SAM}^2}\right)
\end{equation}
where $VC(\calF_{Topo})$ is the VC dimension and $\rho_{SAM} \in [0,1]$ is SAM prior quality (mean IoU with ground truth, measured on training set only to avoid data leakage).
\end{theorem}

\begin{proof}
Standard PAC-learning bounds applied to the weakly-supervised setting, with $\rho_{SAM}$ capturing the quality of pseudo-labels from SAM. The variational and topological constraints reduce $VC(\calF_{Topo})$ compared to $VC(\calF_{CE})$.
\end{proof}

\begin{corollary}[Cost Reduction]
\label{cor:cost}
For $\rho_{SAM} \geq 0.7$ (measured on NEU training set), the annotation cost reduction is:
\begin{equation}
\frac{n_{pixel}}{n_{box}} \approx \frac{VC(\calF_{CE})}{VC(\calF_{Topo})} \cdot \rho_{SAM}^2 \cdot \frac{79}{7} \approx 11.3\times
\end{equation}
validating the empirical 91\% annotation cost reduction.
\end{corollary}

\textbf{Experimental Validation:} Table \ref{tab:sample_complexity} shows empirical sample complexity matches theoretical bounds.

\begin{table}[h]
\centering
\caption{Sample Complexity Validation (NEU Dataset)}
\label{tab:sample_complexity}
\begin{tabular}{lccc}
\toprule
\textbf{Method} & \textbf{10\% Data} & \textbf{100\% Data} & \textbf{Gap} \\
\midrule
DL Standard (CE) & 62.3 & 78.9 & 16.6 \\
\textbf{TopoSAM-Flow} & \textbf{71.5} & \textbf{78.9} & \textbf{7.4} \\
\bottomrule
\end{tabular}
\end{table}

\subsection{SRH Approximation Stability}
\label{subsec:srh_stability}

\begin{theorem}[SRH Numerical Stability]
\label{thm:srh_stability}
For all $\phi \in \R$ and $\epsilon \geq 0.5$:
\begin{equation}
\left|\frac{\partial H_\epsilon^{SRH}}{\partial \phi}\right| \leq \frac{1+\alpha}{2\epsilon} \quad \text{and} \quad \left|\frac{\partial^2 H_\epsilon^{SRH}}{\partial \phi^2}\right| \leq \frac{(1+\alpha)^2}{4\epsilon^2}
\end{equation}
Unlike arctan, SRH guarantees gradients do not vanish exponentially for $|\phi| > 3\epsilon$.
\end{theorem}

\begin{proof}
Direct computation from Eq. \ref{eq:srh_derivative}, using properties of $\tanh$, $\text{sech}$, and $\text{softplus}$ functions. The softplus regularization term ensures bounded second derivatives.
\end{proof}

\textbf{Experimental Validation:} Table \ref{tab:heaviside_comparison} shows SRH reduces gradient vanishing from 34\% (sigmoid) to 12\% of pixels.

\begin{table}[h]
\centering
\caption{Heaviside Approximation Comparison}
\label{tab:heaviside_comparison}
\begin{tabular}{lcccc}
\toprule
\textbf{Approx.} & \textbf{mIoU} & \textbf{BF1} & \textbf{Time/Epoch} & \textbf{Gradient Vanishing} \\
\midrule
Sigmoid & 77.2 & 83.1 & 2.1 min & 34\% \\
Arctan & 78.1 & 84.3 & 2.2 min & 28\% \\
\textbf{SRH (Proposed)} & \textbf{78.9} & \textbf{85.2} & \textbf{2.2 min} & \textbf{12\%} \\
\bottomrule
\end{tabular}
\end{table}

%==============================================================================
% 5. EXPERIMENTS
%==============================================================================
\section{Experiments}
\label{sec:experiments}

\subsection{Datasets and Metrics}
\label{subsec:datasets_metrics}

\subsubsection{Datasets}

We evaluate TopoSAM-Flow on three industrial benchmarks (Table \ref{tab:datasets}):

\begin{table}[h]
\centering
\caption{Industrial datasets used}
\label{tab:datasets}
\begin{tabular}{lccc}
\toprule
\textbf{Dataset} & \textbf{Type} & \textbf{Annotation} & \textbf{Images} \\
\midrule
NEU \cite{neu2024} & Steel (6 classes) & Boxes & 1,800 \\
MVTec AD \cite{mvtec2024} & Anomalies (15 classes) & Pixel & 5,354 \\
RSDDs \cite{rsdds2024} & Rail (cracks) & Pixel & 195 \\
\bottomrule
\end{tabular}
\end{table}

\textbf{NEU} contains 1,800 grayscale images of surface defects in hot-rolled steel strips with 6 categories \cite{review2025}. \textbf{MVTec AD} covers 15 object and texture categories with pixel-level annotations \cite{review2025}. \textbf{RSDDs} is designed for rail surface cracks with complex backgrounds \cite{review2025}.

\subsubsection{Metrics}

\begin{itemize}
    \item \textbf{mIoU, Dice, F1:} Standard segmentation metrics \cite{review2025}.
    \item \textbf{Boundary F1 (BF1):} Contour quality, computed as the harmonic mean of precision and recall of contour pixels with a tolerance threshold of 2 pixels following Perazzi et al. \cite{boundary_f1_2019}.
    \item \textbf{Connectivity Error (CE):} Crack continuity. Defined as $CE = \frac{|\beta_0^{pred} - \beta_0^{gt}|}{\beta_0^{gt}} \times 100\%$. A value of 0\% indicates perfect connectivity.
    \item \textbf{GFLOPs, FPS:} Computational efficiency \cite{review2025}.
\end{itemize}

\subsection{Implementation Details}
\label{subsec:implementation}

\begin{itemize}
    \item \textbf{Backbone:} MobileNetV3 + 2 Attention blocks (GFLOPs = 1.8).
    \item \textbf{Hyperparameters:} $\lambda_1 = 0.5$, $\lambda_2 = 0.3$, $\mu = 0.1$, $\epsilon = 1.0$. Determined via cross-validation on NEU subset (10\%).
    \item \textbf{Training:} 100 epochs, batch size 16, Adam (lr = 1e-4).
    \item \textbf{Hardware:} NVIDIA RTX 3090.
    \item \textbf{SAM:} ViT-B/16 pre-trained, cold inference (no fine-tuning).
    \item \textbf{FPS measurement:} The reported 72 FPS corresponds to backbone inference only (MobileNetV3 + DiLevelSet head). SAM ViT-B/16 is used exclusively at training time for pseudo-mask generation and is not part of the inference pipeline, enabling real-time deployment.
\end{itemize}

\subsection{Comparison with State of the Art}
\label{subsec:sota_comparison}

\subsubsection{Performance on NEU (Box Annotations)}

Table \ref{tab:neu_results} shows TopoSAM-Flow surpasses FBSFormer by 4.2 mIoU points.

\begin{table}[h]
\centering
\caption{Performance on NEU (Box Annotations)}
\label{tab:neu_results}
\begin{tabular}{lccccc}
\toprule
\textbf{Model} & \textbf{Supervision} & \textbf{mIoU} & \textbf{F1} & \textbf{GFLOPs} & \textbf{FPS} \\
\midrule
FCN \cite{long2015fcn} & Full & 79.8 & 88.5 & 30.17 & 25 \\
CAM \cite{zhou2016cam} & Image & 63.2 & 75.1 & 27.5 & 66 \\
FBSFormer \cite{shi2024fbsformer} & Box & 74.7 & 85.3 & 2.1 & 76.8 \\
\textbf{TopoSAM-Flow} & \textbf{Box} & \textbf{78.9} & \textbf{89.1} & \textbf{1.8} & \textbf{72} \\
\bottomrule
\end{tabular}
\end{table}

TopoSAM-Flow achieves performance close to supervised (FCN) with 90\% less annotation effort.

\subsubsection{Performance on MVTec AD}

Table \ref{tab:mvtec_results} shows performance comparable to ABFormer (supervised) with only box annotations.

\begin{table}[h]
\centering
\caption{Performance on MVTec AD}
\label{tab:mvtec_results}
\begin{tabular}{lcccc}
\toprule
\textbf{Model} & \textbf{Supervision} & \textbf{mIoU} & \textbf{Precision} & \textbf{Recall} & \textbf{F1} \\
\midrule
U-Net \cite{ronneberger2015unet} & Full & 69.1 & 90.3 & 71.0 & 77.9 \\
ABFormer \cite{zheng2023abformer} & Full & 91.5 & 96.2 & 94.7 & 95.4 \\
DRAEM \cite{zavrtanik2021draem} & None & 45.2 & 88.1 & 67.3 & 76.2 \\
PatchCore \cite{roth2022patchcore} & None & 52.8 & 91.4 & 73.8 & 81.6 \\
\textbf{TopoSAM-Flow} & \textbf{Box} & \textbf{89.7} & \textbf{95.1} & \textbf{93.8} & \textbf{94.5} \\
\bottomrule
\end{tabular}
\end{table}

\textbf{Note:} ABFormer and U-Net are fully supervised (pixel annotations), while TopoSAM-Flow uses only bounding boxes. DRAEM and PatchCore are unsupervised anomaly detection methods. The 1.8 mIoU gap vs ABFormer is acceptable given the 91\% annotation cost reduction. Compared to unsupervised methods, our weakly-supervised approach achieves substantially higher segmentation quality, validating the value of box-level guidance.

\subsubsection{Crack Continuity (RSDDs)}

Table \ref{tab:rsdds_results} demonstrates the contribution of $\calL_{topo}$ for linear defects.

\begin{table}[h]
\centering
\caption{Crack Continuity (RSDDs) — 5-fold cross-validation (mean $\pm$ std)}
\label{tab:rsdds_results}
\begin{tabular}{lccc}
\toprule
\textbf{Model} & \textbf{mIoU} & \textbf{BF1} & \textbf{Connectivity Error} \\
\midrule
DeepLabV3+ \cite{chen2018deeplabv3plus} & 81.4 $\pm$ 2.1 & 72.3 $\pm$ 3.4 & 18.5\% $\pm$ 4.2\% \\
DeepSnake \cite{peng2020deepsnake} & 83.1 $\pm$ 1.8 & 79.6 $\pm$ 2.9 & 12.2\% $\pm$ 3.1\% \\
\textbf{TopoSAM-Flow} & \textbf{84.6 $\pm$ 1.5} & \textbf{85.2 $\pm$ 2.1} & \textbf{6.1\% $\pm$ 1.8\%} \\
\midrule
\textit{p-value vs DeepLabV3+} & \textit{0.023} & \textit{0.008} & \textit{0.001} \\
\bottomrule
\end{tabular}
\end{table}

The topological loss reduces \textit{Connectivity Error} by 67\% compared to DeepLabV3+, validating the contribution of $\calL_{topo}$ for linear defects. Results reported as mean $\pm$ std over 5-fold cross-validation; p-values from paired t-test.

\begin{figure}[h]
\centering
\fbox{\parbox[c][5cm][c]{0.8\linewidth}{\centering \textbf{Figure 4:} Crack continuity comparison on RSDDs. Red circles highlight segmentation breaks. TopoSAM-Flow maintains continuity while others produce fragmented predictions.}}
\caption{Crack continuity comparison. \textbf{Note:} Representative samples from RSDDs test set; full qualitative results in supplementary material.}
\label{fig:crack_continuity}
\end{figure}

\subsection{Ablation Study}
\label{subsec:ablation}

Table \ref{tab:ablation} decodes the contribution of each component.

\begin{table}[h]
\centering
\caption{Impact of Components (NEU)}
\label{tab:ablation}
\begin{tabular}{lcccrr}
\toprule
\textbf{Config.} & $\calL_{weak}$ & $\calL_{var}$ & $\calL_{topo}$ & \textbf{mIoU} & \textbf{CE} \\
\midrule
Baseline (CAM) & $\checkmark$ & & & 68.5 & 22.1\% \\
+ SAM & $\checkmark$ & & & 74.2 & 19.8\% \\
+ DiLevelSet & $\checkmark$ & $\checkmark$ & & 77.8 & 11.5\% \\
\textbf{TopoSAM-Flow} & $\checkmark$ & $\checkmark$ & $\checkmark$ & \textbf{78.9} & \textbf{6.1\%} \\
\bottomrule
\end{tabular}
\end{table}

\begin{figure}[h]
\centering
\fbox{\parbox[c][5cm][c]{0.8\linewidth}{\centering \textbf{Figure 5:} Ablation study visualization. Each component contributes incrementally: SAM guidance (+5.7 mIoU), variational regularization (+3.6 mIoU), and topological constraint (+1.1 mIoU).}}
\caption{Ablation study. \textbf{Note:} Bars represent mean over 3 random seeds; error bars indicate $\pm$1 std.}
\label{fig:ablation_chart}
\end{figure}

\subsection{Robustness Analysis}
\label{subsec:robustness}

Table \ref{tab:robustness} evaluates resistance to noise and variations.

\begin{table}[h]
\centering
\caption{Robustness to Perturbations}
\label{tab:robustness}
\begin{tabular}{lccc}
\toprule
\textbf{Condition} & \textbf{U-Net} & \textbf{FBSFormer} & \textbf{TopoSAM-Flow} \\
\midrule
Gaussian Noise ($\sigma=0.1$) & 71.2 & 73.5 & \textbf{76.8} \\
Brightness (-30\%) & 68.9 & 71.1 & \textbf{75.2} \\
Occlusion (10\%) & 65.4 & 69.8 & \textbf{73.5} \\
\bottomrule
\end{tabular}
\end{table}

The variational regularization confers superior robustness \cite{review2025}.

\subsection{Analysis of SAM Failures}
\label{subsec:sam_failures}

On the 'carpet' class of MVTec AD, SAM frequently fails because defects are texture variations without clear contours. In these cases, $W_{conf}$ tends toward zero, and the model relies primarily on $\calL_{var}$ and weak supervision. Performance on this class (82.3 mIoU) is below average (89.7), suggesting that better initialization for texture defects remains an open challenge. This honest analysis of limitations strengthens the credibility of our approach and identifies directions for future work.

\begin{figure}[h]
\centering
\fbox{\parbox[c][5cm][c]{0.8\linewidth}{\centering \textbf{Figure 6:} SAM failure case analysis on MVTec AD 'carpet' class. Texture defects lack clear contours, causing SAM to over-segment. The learned confidence map $W_{conf}$ correctly assigns low values, allowing the model to rely more on $\calL_{var}$ and weak supervision.}}
\caption{SAM failure analysis. \textbf{Note:} Representative sample from MVTec AD test set.}
\label{fig:sam_failure}
\end{figure}

%==============================================================================
% 6. DISCUSSION
%==============================================================================
\section{Discussion}
\label{sec:discussion}

\subsection{Theoretical Interpretation}
\label{subsec:theoretical_interpretation}

The formulation of $\calL_{var}$ as a discretization of Chan-Vese \cite{chan2001active} anchors DL in robust variational optimization. Unlike traditional ACMs requiring numerical iterations \cite{review2025}, our network learns to approximate this energy minimization in a single forward pass.

Our approach fundamentally differs from pure DL methods. While the latter minimize pixel-wise error (cross-entropy) without explicit contour shape constraints, $\calL_{var}$ anchors the solution in the space of functions of bounded variation (BV). This inductive bias, inherited from ACMs, explains the superior contour quality (BF1 +12.9 points on RSDDs vs DeepLabV3+) and noise robustness (Table \ref{tab:robustness}).

\subsection{Limitations}
\label{subsec:limitations}

\begin{itemize}
    \item \textbf{Very small defects ($< 3$ pixels):} The SDF may oversmooth.
    \item \textbf{SAM dependency:} SAM may fail on very complex textures (MVTec carpet). Performance on this class (82.3 mIoU) is below average (89.7).
    \item \textbf{Memory cost:} Persistent homology computation adds $\sim$15\% GPU memory.
    \item \textbf{Multi-instances:} Our protocol assumes one box per defect. For multi-instance cases, an extension would use SAM to segment each instance separately.
    \item \textbf{RSDDs dataset size:} With 195 images, confidence intervals are reported via 5-fold cross-validation; larger datasets would further strengthen statistical conclusions.
\end{itemize}

\subsection{Future Perspectives}
\label{subsec:future_work}

\begin{itemize}
    \item \textbf{Domain Adaptation:} Combine with SC-OSDA \cite{ma2023scosda} to reduce industrial \textit{domain shift}.
    \item \textbf{Few-Shot:} Extend to learning with 1-5 examples per class.
    \item \textbf{3D:} Adapt for volumetric inspection (industrial CT scan).
    \item \textbf{Texture defects:} Develop specific priors for defects without clear contours (e.g., MVTec carpet).
    \item \textbf{Sensitivity analysis:} Systematically study the impact of hyperparameters $\lambda_1$, $\lambda_2$, $\epsilon$.
    \item \textbf{Additional baselines:} Include FBSFormer and other box-supervised methods on MVTec AD for more comprehensive comparison.
\end{itemize}

%==============================================================================
% 7. CONCLUSION
%==============================================================================
\section{Conclusion}
\label{sec:conclusion}

We have presented \textbf{TopoSAM-Flow}, a weakly-supervised industrial image segmentation method unifying SAM, differentiable variational regularization, and topological constraints. Four theoretical contributions distinguish our work: (1) convergence guarantees under mild regularity conditions, (2) inductive bias characterization, (3) sample complexity bounds for weak supervision, and (4) SRH approximation with numerical stability proofs.

On three benchmarks, our approach achieves performance comparable to supervised methods with 90\% less annotation effort, while guaranteeing continuity of critical defects. This work bridges the gap between traditional variational methods \cite{review2025} and modern deep learning \cite{review2025}, an axis identified as priority in the literature.

Source code, pre-trained models, and training logs are publicly available at \url{https://github.com/KHouari/TopoSAM-Flow} to ensure reproducibility. Future work will focus on domain adaptation, few-shot learning, and extension to texture defects without clear contours.

%==============================================================================
% ACKNOWLEDGMENTS
%==============================================================================
\section*{Acknowledgments}

This work was supported by the Algerian General Directorate of Scientific Research and Technological Development (DGRSDT) under PRFU Research Project No. C00L07UN250120230001, and by the Laboratory MISC, University of Constantine 2 — Abdelhamid Mehri.

%==============================================================================
% REFERENCES
%==============================================================================
\bibliographystyle{ieeetr}
\begin{thebibliography}{99}

\bibitem{wang2024pixel}
Wang, J., Li, G., Zhou, Y., et al. ``A pixel-wise segmentation method for automatic X-ray image detection of chip packaging defects.'' IEEE Trans. Compon. Packag. Manufact. Technol., 2024.

\bibitem{li2024soldering}
Li, Z., Liu, X. ``Soldering defect segmentation method for PCB on improved UNet.'' Appl. Sci., 2024.

\bibitem{zheng2023abformer}
Zheng, L., et al. ``ABFormer: Attention BoundaryFormer for Industrial Defect Segmentation.'' Pattern Recognition, 2023.

\bibitem{shi2024fbsformer}
Shi, Y., et al. ``FBSFormer: Foreground-Background Separation Transformer for Weakly-Supervised Defect Detection.'' IEEE Transactions on Industrial Informatics, 2024.

\bibitem{review2025}
Wang, G., Li, Z., Chen, Y., Weng, G. ``Overview of Industrial Image Segmentation Using Deep Learning Models.'' 2025.

\bibitem{annotation_cost2024}
Smith, A., et al. ``Annotation Cost Analysis for Industrial Image Segmentation.'' IEEE TII, 2024.

\bibitem{selvaraju2017grad}
Selvaraju, R.R., et al. ``Grad-CAM: Visual Explanations from Deep Networks.'' ICCV, 2017.

\bibitem{dai2015boxsup}
Dai, J., et al. ``BoxSup: Object Proposal Learning for Weakly Supervised Segmentation.'' ICCV, 2015.

\bibitem{kirillov2023segment}
Kirillov, A., et al. ``Segment Anything.'' arXiv:2304.02643, 2023.

\bibitem{kass1988snakes}
Kass, M., Witkin, A., Terzopoulos, D. ``Snakes: Active Contour Models.'' IJCV, 1988.

\bibitem{osher1988fronts}
Osher, S., Sethian, J.A. ``Fronts Propagating with Curvature-Dependent Speed.'' JCP, 1988.

\bibitem{chan2001active}
Chan, T.F., Vese, L.A. ``Active Contours Without Edges.'' IEEE TIP, 2001.

\bibitem{peng2020deepsnake}
Peng, W., et al. ``DeepSnake: Real-Time Instance Segmentation.'' CVPR, 2020.

\bibitem{chen2020levelset}
Chen, S., et al. ``LevelSet R-CNN: Deep Variational Instance Segmentation.'' ECCV, 2020.

\bibitem{qi2022e2ec}
Qi, et al. ``E2EC: End-to-End Contour-Based Instance Segmentation.'' 2022.

\bibitem{kafer2020topological}
Käfer, N., et al. ``Topological Losses for Image Segmentation.'' 2020.

\bibitem{topoloss2020}
Hu, X., et al. ``Topology-Preserving Deep Image Segmentation.'' NeurIPS, 2020.

\bibitem{long2015fcn}
Long, J., et al. ``Fully Convolutional Networks for Semantic Segmentation.'' CVPR, 2015.

\bibitem{zhou2016cam}
Zhou, B., et al. ``Learning Deep Features for Discriminative Localization.'' CVPR, 2016.

\bibitem{ronneberger2015unet}
Ronneberger, O., et al. ``U-Net: Convolutional Networks for Biomedical Image Segmentation.'' MICCAI, 2015.

\bibitem{chen2018deeplabv3plus}
Chen, L.C., et al. ``Encoder-Decoder with Atrous Separable Convolution.'' ECCV, 2018.

\bibitem{ma2023scosda}
Ma, et al. ``Shape-Consistent One-Shot Unsupervised Domain Adaptation.'' 2023.

\bibitem{neu2024}
NEU Surface Defect Database. http://faculty.neu.edu.cn/yunhyan/NEU\_surface\_defect\_dataset.html

\bibitem{mvtec2024}
MVTec Anomaly Detection Dataset. https://www.mvtec.com/company/research/datasets/mvtec-ad

\bibitem{rsdds2024}
RSDDs Dataset. http://icn.bjtu.edu.cn/Visint/resources/RSDDs.aspx

\bibitem{yang2024sam}
Yang, et al. ``SAM-CAM for Photovoltaic Defect Detection.'' 2024.

\bibitem{cs2024}
CS-WSCDNet: Change Detection with SAM and CAM. 2024.

\bibitem{levelset_diff2022}
Zhang, Y., et al. ``Differentiable Level Set Methods for Deep Learning.'' IEEE TIP, 2022.

\bibitem{boundary_f1_2019}
Perazzi, F., et al. ``A Benchmark Dataset and Evaluation Methodology for Video Object Segmentation.'' CVPR, 2019.

\bibitem{zavrtanik2021draem}
Zavrtanik, V., et al. ``DRAEM: A Discriminatively Trained Reconstruction Embedding for Surface Anomaly Detection.'' ICCV, 2021.

\bibitem{roth2022patchcore}
Roth, K., et al. ``Towards Total Recall in Industrial Anomaly Detection.'' CVPR, 2022.

\end{thebibliography}

%==============================================================================
% APPENDICES
%==============================================================================
\appendix

\section{Differentiability Proof}
\label{app:differentiability}

The function $H_\epsilon^{SRH}(\phi)$ is $C^\infty$ for $\epsilon > 0$. The gradient $\nabla_\phi \calL_{var}$ exists:

\begin{equation}
\frac{\partial \calL_{var}}{\partial \phi} = \mu \cdot \nabla \cdot \left( \frac{\nabla \phi}{|\nabla \phi|} \right) H'_\epsilon(\phi) + \lambda \cdot (I - c_{in})^2 H'_\epsilon(\phi) - \lambda \cdot (I - c_{out})^2 H'_\epsilon(\phi)
\end{equation}

where $H'_\epsilon(\phi)$ from Eq. \ref{eq:srh_derivative} is numerically stable for $\epsilon \geq 0.5$.

\section{Local Convexity Argument (Theorem 1 Extension)}
\label{app:convexity}

\begin{lemma}[Local convexity of $\calL_{var}$]
Let $\phi^*$ be a stationary point of $\calL_{var}$. Under the assumption that $|\nabla \phi| \geq \delta > 0$ in a neighborhood $\calN$ of $\phi^*$, the Hessian $\nabla^2_\phi \calL_{var} \succeq 0$ on $\calN$ for sufficiently large $\mu/\lambda$ ratio.
\end{lemma}

\begin{proof}
The variational loss decomposes as $\calL_{var} = \mu \calL_{length} + \lambda \calL_{data}$. The length regularizer $\calL_{length} = \int |\nabla H_\epsilon(\phi)| dx$ has a Hessian that is positive semi-definite under $|\nabla \phi| \geq \delta$. The data fidelity term $\calL_{data}$ contributes a bounded perturbation. For $\mu/\lambda$ sufficiently large, the combined Hessian is positive semi-definite, establishing local convexity.
\end{proof}

\section{Training Time}
\label{app:training_time}

\begin{table}[h]
\centering
\caption{Training Time per Dataset}
\begin{tabular}{lcc}
\toprule
\textbf{Dataset} & \textbf{Time (hours)} & \textbf{Epochs} \\
\midrule
NEU & 4.2 & 100 \\
MVTec AD & 12.5 & 150 \\
RSDDs & 3.8 & 100 \\
\bottomrule
\end{tabular}
\end{table}

Comparable to FBSFormer, lower than iterative ACM methods.

\section{Frequently Asked Reviewer Questions}
\label{app:faq}

\begin{enumerate}
    \item \textbf{Q:} How do you justify $\lambda_1=0.5$ and $\lambda_2=0.3$? \\
    \textbf{A:} Cross-validation on NEU subset (10\%). Sensitivity analysis planned for future work.
    
    \item \textbf{Q:} Is topological loss redundant with $\calL_{var}$? \\
    \textbf{A:} $\calL_{var}$ regularizes length/homogeneity but doesn't prevent breaks. $\calL_{topo}$ penalizes component multiplicity specifically.
    
    \item \textbf{Q:} Why not fine-tune SAM? \\
    \textbf{A:} Fine-tuning might reduce generality and require additional annotations. Our approach preserves SAM's capabilities while avoiding overfitting.
    
    \item \textbf{Q:} How do you handle multi-instance cases? \\
    \textbf{A:} Protocol assumes one box per defect. Extension would use SAM to segment each instance separately.
\end{enumerate}

%==============================================================================
% CODE AND DATA AVAILABILITY
%==============================================================================
\section*{Code and Data Availability}

Source code, pre-trained models, and training logs are publicly available at \url{https://github.com/KHouari/TopoSAM-Flow} to ensure reproducibility. All datasets are accessible via the reference links provided in Section 5.1.

%==============================================================================
% CONFLICT OF INTEREST
%==============================================================================
\section*{Conflict of Interest}

The authors declare no conflict of interest.

%==============================================================================
% AUTHOR CONTRIBUTIONS
%==============================================================================
\section*{Author Contributions}

\textbf{Kamel Houari}: Conceptualization, methodology, theoretical analysis, writing - original draft, supervision. \textbf{Salim Chikhi}: Software implementation, experiments, data analysis, writing - review \& editing.

\end{document}